\documentclass[../../book.tex]{subfiles}
\begin{document}
\graphicspath{{images/}}
\chaptertitle{Lab 1: USV Modeling and System Identification}{3}

\section{Learning Objectives}

By the end of this lab you should be able to:

\begin{description}
  \item[\textbf{Derive}] Linear transfer function models for a marine vessel from first principles (Newton's second law).
  \item[\textbf{Identify}] The physical meaning of the DC gain and time constant of a first-order plant.
  \item[\textbf{Describe}] How propeller thrust controls surge speed and how rudder moment controls heading.
  \item[\textbf{Simulate}] The \textbf{open-loop} step response for nominal USV parameters, to predict the dynamics of speed and yaw response to manual inputs.
  \item[\textbf{Collect}] Collect experimental step-response data from a USV running ArduPilot and plot the open-loop response, measuring step-response metrics: rise time, settling time, and percent overshoot.
  \item[\textbf{System Identification}] Fit the model to the data by selecting transfer function parameters so that the model output matches the measured step response.
\end{description}


\section{Deliverable}

You will work in small teams on this lab.  Each team will create a short document with the following elements.  Each member of the team will submit a copy of the document (with all the team members names) through Sakai.

\begin{itemize}
  \item Title and names of team members.
  \item Nominal Response: Figures for the open-loop step response for surge and yaw-rate using the provided nominal parameters.  Include a narrative that compares the expected response metrics (rise-time, overshoot, steady-state error) with the observed metrics.
  \item Experimental Summary: Figures of the experimental results for surge and yaw-rate.   This is soley based on the data.  Include a narrative to identify the empirical system parameters for each response: time-constant and DC gain.
  \item System Identification: After completing the system identification (via curve fitting), present figures with experimental open-loop response (the data) and the identified transfer function response (the model) on the same axes.  Include a narrative to quantitatively compare the model to the experiment with metrics such as rise time, settling time, and percent overshoot.  Include a table for a side-by-side comparision of experimental and model metrics.
\end{itemize}

The \emph{narratives} should be brief to explain the key associated figures and metrics.  The narratives should not include the process, just the results.

% ============================================================
\section{Physical System}

\subsection{Vehicle Description}

The vehicle is a small, single-hull unmanned surface vehicle (USV)
approximately 1.1\,m in length and 5\,kg in mass.  It carries one
fixed-pitch propeller on the centerline for thrust and one stern rudder
for steering.  Because the propeller is on the centerline and because we are using the model for control, we can consider thrust and
steering to be effectively decoupled --- thrust controls forward speed;
the rudder controls heading direction.

\subsection{Coordinate Frame and Degrees of Freedom}

We adopt the standard SNAME body-fixed frame.  The $x$-axis points
forward (bow), the $y$-axis points to port, and the $z$-axis points
upward.  For this lab we model two scalar motions:

\vspace{6pt}
\begin{center}
\begin{tabular}{C{1.4cm} C{1.4cm} C{1.4cm} L{8.2cm}}
  \textbf{DOF} & \textbf{Variable} & \textbf{Units} & \textbf{Description} \\
  \midrule
  Surge & $u$ & m/s &
    Forward (longitudinal) velocity along body $x$-axis \\
  Yaw   & $r,\,\psi$ & rad/s, rad &
    Yaw rate $r = \dot{\psi}$; heading angle $\psi$ measured from North \\
  \bottomrule
\end{tabular}
\end{center}
\vspace{6pt}

Sway (lateral velocity) and roll are neglected.  For a slow,
narrow-beam hull this is a reasonable simplification.  The two channels
(surge and yaw) are treated as independent and are controlled by
separate PID loops.

% ============================================================
\section{Model Derivation}

\subsection{Surge Channel}

\subsubsection{Equation of Motion}

Applying Newton's second law in the surge (forward) direction:

\begin{equation}
  m_{\mathrm{eff}}\,\dot{u} \;=\; F_{\mathrm{thrust}} - F_{\mathrm{drag}}
  \label{eq:surge_newton}
\end{equation}

where $m_{\mathrm{eff}}$ is an effective inertia that includes the vessel's
physical mass $m$ plus an added-mass correction $m_a$ due to the entrained
water that accelerates with the hull:
\[
  m_{\mathrm{eff}} = m + m_a \approx 1.20\,m
\]

For the drag force we adopt a \emph{linear} (viscous) model.  In reality,
hydrodynamic drag is approximately quadratic in velocity
($F_{\mathrm{drag}} \propto u^2$), but linearizing about a nominal operating
speed $u_0$ gives:
\[
  F_{\mathrm{drag}} = b_u\,u
\]
where $b_u$ is the linearized surge drag coefficient [N$\cdot$s/m].
Absorbing the 1.20 factor into the lumped parameter $m$, the complete
surge equation of motion is:

\begin{equation}
  \boxed{m\,\dot{u} + b_u\,u = F}
  \label{eq:surge_eom}
\end{equation}

\subsubsection{Transfer Function}

Taking the Laplace transform of \eqref{eq:surge_eom} with zero initial
conditions:
\[
  m\,s\,U(s) + b_u\,U(s) = F(s)
  \qquad\Longrightarrow\qquad
  U(s)\,(ms + b_u) = F(s)
\]

The surge plant transfer function --- from thrust $F$ [N] to surge speed
$U$ [m/s] --- is:

\begin{equation}
  G_{\mathrm{surge}}(s)
  \;=\; \frac{U(s)}{F(s)}
  \;=\; \frac{1}{ms + b_u}
  \;=\; \frac{K_u}{\tau_u\,s + 1}
  \label{eq:Gsurge}
\end{equation}

This is a \textbf{first-order system} with:

\begin{itemize}
  \item \textbf{DC gain:} \;$K_u = 1/b_u$ \;[m\,s$^{-1}$/N]
        --- steady-state speed per unit thrust
  \item \textbf{Time constant:} \;$\tau_u = m/b_u$ \;[s]
        --- time to reach 63\% of final speed
  \item \textbf{Pole:} \;$s = -b_u/m$ \;(left-half plane --- open-loop stable)
\end{itemize}

\subsubsection{Command-Scaled Transfer Function}

In practice the propeller is driven by a speed controller that accepts a
thrust command $\delta_T \in [0,\,100]$ expressed as a percentage of full
thrust.  Assuming a linear actuator model:
\begin{equation}
  F \;=\; F_{\max}\,\frac{\delta_T}{100}
  \label{eq:thrust_cmd}
\end{equation}
where $F_{\max}$ is the maximum thrust produced at 100\% command.
Substituting \eqref{eq:thrust_cmd} into \eqref{eq:surge_eom} and taking the
Laplace transform gives the command-scaled surge transfer function:

\begin{equation}
  G_{\mathrm{surge}}^*(s)
  \;\equiv\; \frac{U(s)}{\Delta_T(s)}
  \;=\; \frac{F_{\max}/100}{ms + b_u}
  \;=\; \frac{K_u^*}{\tau_u\,s + 1}
  \label{eq:Gsurge_cmd}
\end{equation}

where the command-scaled DC gain and time constant are:
\[
  K_u^* \;=\; \frac{F_{\max}}{100\,b_u}
  \quad \left[\frac{\mathrm{m/s}}{\%}\right],
  \qquad
  \tau_u \;=\; \frac{m}{b_u} \quad [\mathrm{s}]
\]

The time constant $\tau_u$ is unchanged --- it depends only on the vehicle
dynamics, not on the actuator scaling.  Only the DC gain absorbs the
actuator gain $F_{\max}/100$.  This means that from a step-response
experiment, $\tau_u$ can be read directly from the response shape while
$K_u^*$ is determined from the steady-state speed achieved at a known
command percentage.

% ─────────────────────────────────────────────
\subsection{Yaw Channel}

\subsubsection{Equation of Motion}

Applying Newton's second law for rotation about the vertical ($z$) axis:

\begin{equation}
  I_{\mathrm{eff}}\,\dot{r} \;=\; N_{\mathrm{rudder}} - N_{\mathrm{drag}}
  \label{eq:yaw_newton}
\end{equation}

where $I_{\mathrm{eff}} = I + I_a$ is the effective yaw inertia (hull
inertia plus rotational added inertia), $r = \dot\psi$ is the yaw rate,
$N_{\mathrm{rudder}}$ is the hydrodynamic moment produced by the rudder
[N$\cdot$m], and $N_{\mathrm{drag}} = b_r\,r$ is a linear damping moment.
The yaw-rate equation is:

\begin{equation}
  \boxed{I\,\dot{r} + b_r\,r = N}
  \label{eq:yaw_eom}
\end{equation}

\subsubsection{Heading Transfer Function}

Because the commanded output is usually the heading angle $\psi$ rather
than yaw rate $r$, and since $r = \dot\psi$, we note that
$\Psi(s) = R(s)/s$.

The yaw-rate plant transfer function (moment $N$ to yaw rate $R$) is the
same first-order form as surge:
\[
  G_r(s) \;=\; \frac{R(s)}{N(s)} \;=\; \frac{1}{Is + b_r}
\]

The heading plant transfer function (moment $N$ to heading angle $\Psi$)
adds an integrator:

\begin{equation}
  G_{\mathrm{yaw}}(s)
  \;=\; \frac{\Psi(s)}{N(s)}
  \;=\; \frac{1}{s\,(Is + b_r)}
  \;=\; \frac{K_r/\tau_r}{s\,(\tau_r\,s + 1)}
  \label{eq:Gyaw}
\end{equation}

  The heading plant \eqref{eq:Gyaw} has a \textbf{free integrator} (a pole
  at $s = 0$).  Physically: a constant rudder moment produces a constantly
  growing heading angle.  As a control-design consequence, even a
  proportional-only controller achieves zero steady-state error to a step
  heading command --- verify this with the Final Value Theorem as part of
  your pre-lab.

\subsubsection{Command-Scaled Transfer Function}

The rudder servo accepts a displacement command $\delta_R \in [-100,\,100]$
expressed as a percentage of full deflection, with positive values
corresponding to a port (left) turn.  Assuming a linear mapping from
command to hydrodynamic moment:
\begin{equation}
  N \;=\; N_{\max}\,\frac{\delta_R}{100}
  \label{eq:rudder_cmd}
\end{equation}
where $N_{\max}$ is the maximum yaw moment at full rudder deflection.
Substituting \eqref{eq:rudder_cmd} into \eqref{eq:yaw_eom} yields the
command-scaled yaw-rate transfer function:

\begin{equation}
  G_r^*(s)
  \;\equiv\; \frac{R(s)}{\Delta_R(s)}
  \;=\; \frac{N_{\max}/100}{Is + b_r}
  \;=\; \frac{K_r^*}{\tau_r\,s + 1}
  \label{eq:Gr_cmd}
\end{equation}

and the command-scaled heading transfer function:

\begin{equation}
  G_{\mathrm{yaw}}^*(s)
  \;\equiv\; \frac{\Psi(s)}{\Delta_R(s)}
  \;=\; \frac{N_{\max}/100}{s\,(Is + b_r)}
  \;=\; \frac{K_r^*}{s\,(\tau_r\,s + 1)}
  \label{eq:Gyaw_cmd}
\end{equation}

where:
\[
  K_r^* \;=\; \frac{N_{\max}}{100\,b_r}
  \quad \left[\frac{\mathrm{rad/s}}{\%}\right],
  \qquad
  \tau_r \;=\; \frac{I}{b_r} \quad [\mathrm{s}]
\]

As with surge, the time constant $\tau_r$ is unaffected by actuator
scaling.  Note that because the heading plant contains a free integrator,
a step rudder command produces a ramp in heading angle at steady state ---
the vehicle turns continuously.  System identification for the yaw channel
is therefore performed on the \emph{yaw-rate} response $r(t)$, which does
reach a steady state, and the heading model is then constructed by
appending the integrator.

% ============================================================
\section{Nominal Parameters}

The following lumped parameters are used in simulation.  They are
representative of a \(\sim\)5\,kg, 1.1\,m recreational-grade USV.

\vspace{6pt}
\begin{center}
\begin{tabular}{C{1.5cm} C{1.5cm} C{2.8cm} L{7.4cm}}
  \textbf{Symbol} & \textbf{Value} & \textbf{Units} & \textbf{Description} \\
  \midrule
  $m$       & 5.0   & kg                    & Effective surge mass (hull + added mass) \\
  $b_u$     & 8.0   & N$\cdot$s/m           & Linear surge drag coefficient \\
  $I$       & 0.8   & kg$\cdot$m$^2$        & Effective yaw inertia (hull + added inertia) \\
  $b_r$     & 1.2   & N$\cdot$m$\cdot$s/rad & Linear yaw damping coefficient \\
  $\tau_u$  & 0.625 & s                     & Surge time constant $\tau_u = m/b_u$ \\
  $\tau_r$  & 0.667 & s                     & Yaw-rate time constant $\tau_r = I/b_r$ \\
  $K_u$     & 0.125 & (m/s)/N               & Surge DC gain (physical input) $K_u = 1/b_u$ \\
  $K_r$     & 0.833 & (rad/s)/(N$\cdot$m)   & Yaw-rate DC gain (physical input) $K_r = 1/b_r$ \\
  $K_u^*$   & ---   & (m/s)/\%              & Surge DC gain (command-scaled) $K_u^* = F_{\max}/(100\,b_u)$ \\
  $K_r^*$   & ---   & (rad/s)/\%            & Yaw-rate DC gain (command-scaled) $K_r^* = N_{\max}/(100\,b_r)$ \\
  \bottomrule
\end{tabular}
\end{center}
\vspace{3pt}
{\small The command-scaled gains $K_u^*$ and $K_r^*$ depend on $F_{\max}$ and
$N_{\max}$, which are vehicle-specific and determined by system identification
(Section~\ref{sec:sysid}).  Dashes indicate that nominal values are not
assigned a priori.}
\vspace{6pt}

% ============================================================
\section{Transfer Function Summary}

\subsection*{Physical-input forms (force/moment in N and N$\cdot$m)}
\begin{center}
\begin{tabular}{L{3.0cm} L{5.5cm} L{4.7cm}}
  \textbf{Name} & \textbf{Expression} & \textbf{Notes} \\
  \midrule
  Surge plant &
    $\displaystyle G_{\mathrm{surge}}(s) = \dfrac{1}{ms + b_u}$ &
    $F\,[\mathrm{N}] \to u\,[\mathrm{m/s}]$; pole at $s=-b_u/m$ \\[10pt]

  Yaw-rate plant &
    $\displaystyle G_r(s) = \dfrac{1}{Is + b_r}$ &
    $N\,[\mathrm{N{\cdot}m}] \to r\,[\mathrm{rad/s}]$; pole at $s=-b_r/I$ \\[10pt]

  Heading plant &
    $\displaystyle G_{\mathrm{yaw}}(s) = \dfrac{1}{s\,(Is + b_r)}$ &
    $N\,[\mathrm{N{\cdot}m}] \to \psi\,[\mathrm{rad}]$; poles at $0,\,-b_r/I$ \\[10pt]
  \bottomrule
\end{tabular}
\end{center}

\subsection*{Command-scaled forms (input in \% of full actuation)}
\begin{center}
\begin{tabular}{L{3.0cm} L{5.5cm} L{4.7cm}}
  \textbf{Name} & \textbf{Expression} & \textbf{Notes} \\
  \midrule
  Surge plant &
    $\displaystyle G_{\mathrm{surge}}^*(s) = \dfrac{K_u^*}{\tau_u\,s + 1}$ &
    $\delta_T\,[\%] \to u\,[\mathrm{m/s}]$;\ \ $K_u^* = \dfrac{F_{\max}}{100\,b_u}$ \\[12pt]

  Yaw-rate plant &
    $\displaystyle G_r^*(s) = \dfrac{K_r^*}{\tau_r\,s + 1}$ &
    $\delta_R\,[\%] \to r\,[\mathrm{rad/s}]$;\ \ $K_r^* = \dfrac{N_{\max}}{100\,b_r}$ \\[12pt]

  Heading plant &
    $\displaystyle G_{\mathrm{yaw}}^*(s) = \dfrac{K_r^*}{s\,(\tau_r\,s + 1)}$ &
    $\delta_R\,[\%] \to \psi\,[\mathrm{rad}]$ \\[10pt]
  \bottomrule
\end{tabular}
\end{center}

The parameters to be identified from experiment are $K_u^*$, $\tau_u$,
$K_r^*$, and $\tau_r$.  The time constants are shared between the physical
and command-scaled forms; only the DC gains differ by the actuator scale
factor.

% ============================================================
\section{System Identification from Step Responses}
\label{sec:sysid}

\subsection{Overview}

System identification is the process of determining model parameters from
measured input--output data.  For a first-order system the step response
has a simple analytical form that makes parameter extraction
straightforward.  The procedure below applies to both the surge channel
(input: $\delta_T$; output: $u$) and the yaw channel (input: $\delta_R$;
output: $r$).  In each case a \emph{step command} is applied --- the input
is held at zero, then switched to a constant value $\delta_0$ and held
there --- and the output is recorded.

\subsection{Analytical Step Response of a First-Order System}

For a generic first-order command-scaled plant
\[
  G^*(s) = \frac{K^*}{\tau s + 1}
\]
the response to a step input of magnitude $\delta_0$ (applied at $t=0$,
starting from rest) is:

\begin{equation}
  y(t) \;=\; K^*\,\delta_0\,\bigl(1 - e^{-t/\tau}\bigr), \qquad t \geq 0
  \label{eq:step_response}
\end{equation}

Two key features of \eqref{eq:step_response} are used for identification:

\begin{itemize}
  \item \textbf{Steady-state value:} as $t \to \infty$,
        $y_\infty = K^*\,\delta_0$.  Therefore:
        \begin{equation}
          K^* \;=\; \frac{y_\infty}{\delta_0}
          \label{eq:K_id}
        \end{equation}

  \item \textbf{Time constant:} at $t = \tau$ the output reaches
        $63.2\%$ of its final value, i.e.\ $y(\tau) = 0.632\,y_\infty$.
        Therefore $\tau$ is the time at which the measured response first
        crosses $0.632\,y_\infty$:
        \begin{equation}
          \tau \;=\; t\big|_{y\,=\,0.632\,y_\infty}
          \label{eq:tau_id}
        \end{equation}
\end{itemize}

A complementary graphical check is the \emph{initial slope}: at $t=0$,
$\dot{y}(0) = K^*\delta_0/\tau$, so a tangent line drawn at the origin
intersects the steady-state value at $t = \tau$.  This provides a
second, independent estimate of $\tau$.

\subsection{Surge Model System Identification}

The ArduPilot autopilot firmware should be configured to log all internal data to the removeable SD card.  To check this, verify in MP that \texttt{LOG\_BITMASK = 65535}.

\begin{enumerate}
  \item Setup: Designate a pilot (driving the USV swith R/C transmitter), a planner (to monitor in mission planner) and a scribe (taking notes with time stamps).  For each experiment trial, have the scribe note the time and any key observations.  This will make it much easer to sort through the data post-experiment. 
  \item Speed response trials: 
  \begin{itemize}
    \item Prior to each trial, the scribe should note the time and record any identifying behavior.
    \item With the vehicle at rest on open water have the pilot apply a step
        thrust command $\delta_T = \delta_0$ (a value in the range 20--50\%
        is recommended to remain in the approximately linear regime while
        producing a measurable response).  Hold the command constant until
        the surge speed $u$ reaches a clear steady state (typically
        $5\tau_u$ or about 10\,s).  When in doubt, hold the command longer to ensure you capture steady-state.
      \item During the experiment, if the planner can record 
        $u(t)$ (forward speed) and R/C command from mission planner, it wil further help with post-processing.
      \item Conduct 2-3 successful straight line surge tests.
  \end{itemize}

\end{enumerate}


\subsection{Yaw Model System Identification}


\begin{enumerate}
  \item With the vehicle moving at a low, approximately constant surge
        speed (so that rudder authority is meaningful), apply a step
        rudder command $\delta_R = \delta_0$ and hold it constant.
        Record the \emph{yaw rate} $r(t)$ from the IMU gyroscope.
\end{enumerate}

\subsubsection{Practical Notes}

\begin{itemize}
  \item Perform the yaw experiment at a fixed, repeatable surge speed.
        Rudder moment depends on flow velocity over the rudder, so
        $N_{\max}$ (and therefore $K_r^*$) will vary with speed.  Note
        the surge speed at which identification was performed; the yaw
        model is only valid near that operating point.
  \item Use symmetric step commands ($+\delta_0$ then $-\delta_0$) and
        average the identified gains to reduce the effect of any
        asymmetry in the rudder or hull.
\end{itemize}

\subsubsection{Verify Logged Data}

ArduPilot logs data in a binary format called \texttt{.bin}.  The files are stored on the SD card and can be downloaded to a computer for analysis.  It is good practice to verify that the data is being logged correctly before leaving the field.  We will discuss methods of doing this during the lab session.

\subsection{System Identification: Fitting the Model to Data}


\subsubsection{Graphical Method}

To get a first rough estimate of the parameters, plot the measured step response $y(t)$ and apply the formulas \eqref{eq:K_id} and \eqref{eq:tau_id} to extract $K^*$ and $\tau$ for each of the surge and yaw models. This is a quick way to get ballpark estimates without needing to do any curve fitting.  We will also use these estimates as initial guesses for the more rigorous curve-fitting method described next.

Produce two separate figures. For the surge channel, plot $u(t)$ and annotate the steady-state value $y_\infty$ and the time at which the response first crosses $0.632\,y_\infty$.  For the yaw channel, plot $r(t)$ and similarly annotate the steady-state value and the 63.2\% crossing time.  On these same axes, also plot the step response predicted by the identified parameters to visually check the fit.

\subsubsection{Curve Fitting}

Once initial estimates $\hat K^*$ and $\hat\tau$ have been obtained
graphically, a better fit can be achieved by minimizing the least-squares
error between the analytical step response \eqref{eq:step_response} and
the measured data.  Given $N$ data samples $(t_k,\,y_k)$, define the
residual vector:
\[
  \varepsilon_k \;=\; y_k \;-\; \hat K^*\,\delta_0\,
  \bigl(1 - e^{-t_k/\hat\tau}\bigr)
\]
and minimize $\sum_k \varepsilon_k^2$ over $(K^*, \tau)$.
In MATLAB this is accomplished with \texttt{lsqcurvefit}, which is inlucled in the Optimization Toolbox.  The code snippet below shows how to set up the optimization problem:

  \begin{verbatim}
  function y = first_order_step(params, t, delta0)
    K_star = params(1);
    tau = params(2);
    y = K_star * delta0 * (1 - exp(-t / tau));
  end

  % t_data, y_data: arrays from experiment; delta0: step magnitude
  params0 = [K_hat, tau_hat];  % initial guesses from graphical method
  params_fit = lsqcurvefit(...
    @(p, t) first_order_step(p, t, delta0), ...
    params0, t_data, y_data);
  K_fit = params_fit(1);
  tau_fit = params_fit(2);
  \end{verbatim}

The graphical estimates serve as initial guesses \texttt{p0}; providing
good starting values is important because the optimization is nonlinear
in $\tau$.

\subsection{Validation}

After fitting always validate the identified model against data not used in the fit.  Create your transfer function models for surge and yaw-rate using the identified parameters from the curve fitting, and plot the model-predicted step response on top of the measured data.  Check that:
\begin{itemize}
  \item The model captures the key features of the response: rise time,
        settling time, and percent overshoot.
  \item Report the root-mean-square error:
        \[
          \mathrm{RMSE} \;=\; \sqrt{\frac{1}{N}\sum_{k=1}^N \varepsilon_k^2}
        \]
        as a quantitative measure of fit quality.
\end{itemize}

\subsection{Summary of Identified Parameters}

Record your identified parameters in the following table:

\vspace{6pt}
\begin{center}
\begin{tabular}{C{2.2cm} C{2.2cm} C{2.2cm} C{2.2cm} C{2.2cm}}
  \textbf{Channel} & \textbf{$\delta_0$ [\%]} & \textbf{$y_\infty$} &
  \textbf{$K^*$ (identified)} & \textbf{$\tau$ (identified)} \\
  \midrule
  Surge & & & & \\[8pt]
  Yaw rate & & & & \\[8pt]
  \bottomrule
\end{tabular}
\end{center}
\vspace{6pt}

\noindent Units: $K_u^*$ in (m/s)/\%;\; $K_r^*$ in (rad/s)/\%;\;
$\tau_u,\,\tau_r$ in s.

\end{document}
